\begin{tabularx}{\textwidth}[t]{|X|X|}
  % \hline \textbf{Notions et contenus} & \textbf{Capacités exigibles} \\
  % \hline \multicolumn{2}{|l|}{ 4.1. Changements de référentiel } \\
  % \hline 
  %   Référentiel en translation rectiligne uniforme 
  %   par rapport à un autre : transformation de 
  %   Galilée, composition des vitesses.
  %  & 
  %   Relier la transformation de Galilée et la 
  %   formule de composition des vitesses à la 
  %   relation de Chasles et au caractère 
  %   supposé absolu du temps.
  % \\
  % \hline 
  %   Composition des vitesses et des 
  %   accélérations dans le cas d'un référentiel en 
  %   translation par rapport à un autre : point 
  %   coïncident, vitesse d'entraînement, 
  %   accélération d'entraînement.
  %  & 
  %   Exprimer la vitesse d'entraînement et 
  %   l'accélération d'entraînement.
  % \\
  % \hline 
  %   Composition des vitesses et des accélérations 
  %   dans le cas d'un référentiel en rotation 
  %   uniforme autour d'un axe fixe : point 
  %   coïncident, vitesse d'entraînement, 
  %   accélération d'entraînement, accélération de 
  %   Coriolis.
  %  & 
  %   Exprimer la vitesse d'entraînement et 
  %   l'accélération d'entraînement. 
  %   Citer et utiliser l'expression de l'accélération 
  %   de Coriolis.
  % \\
  % \hline

  \hline \textbf{Notions et contenus} & \textbf{Capacités exigibles} \\
  \hline  \multicolumn{2}{|l|}{ 4.2. Dynamique dans un référentiel non galiléen} \\
  \hline 
    Cas d'un référentiel en translation par 
    rapport à un référentiel galiléen : force 
    d'inertie d'entrainement.
   & 
    Déterminer la force d'inertie d'entraînement. 
    Appliquer la deuxième loi de Newton, le 
    théorème du moment cinétique et le théorème 
    de l'énergie cinétique dans un référentiel non 
    galiléen.
  \\
  \hline 
    Cas d'un référentiel en rotation uniforme 
    autour d'un axe fixe dans un référentiel 
    galiléen : force d'inertie d'entraînement, 
    force d'inertie de Coriolis.
   & 
    Exprimer la force d'inertie d'entraînement et la 
    force d'inertie de Coriolis. 
    Associer la force d'inertie d'entraînement 
    axifuge à l'expression familière «force 
    centrifuge ». 
    Appliquer la deuxième loi de Newton, le 
    théorème du moment cinétique et le théorème 
    de l'énergie cinétique dans un référentiel non 
    galiléen.
  \\
  \hline
\end{tabularx}
